% ch6_b (书页 251-262 / p266-p277)
%--A-TAIL--
%JOIN% 于是有
% ---- 书页 251 / p266 ----
\begin{equation*}
E_{R} = \frac{1}{4\pi G}\int_{\Sigma} d^{3}x\,\sqrt{\gamma}\, n_{\mu}\nabla_{\nu}\left(\nabla^{\mu}K^{\nu}\right).
\tag{6.37}
\end{equation*}

注意，由 Killing 方程升起指标可得 $\nabla^{\mu}K^{\nu}=-\nabla^{\nu}K^{\mu}$。因此，正如前面处理电荷时那样，我们可以再次利用 Stokes 定理，把 $E_{R}$ 写成空间无穷远处的积分，
\begin{equation*}
\boxed{\,E_{R} = \frac{1}{4\pi G}\int_{\partial\Sigma} d^{2}x\,\sqrt{\gamma^{(2)}}\, n_{\mu}\sigma_{\nu}\nabla^{\mu}K^{\nu}.\,}
\tag{6.38}
\end{equation*}
这个表达式就是与类时 Killing 矢量 $K^{\mu}$ 相联系的\textbf{Komar 积分}（Komar integral）；它可以解释为稳态时空的总能量。

为了让我们自己确信方向没错，就用度规为式 (6.18) 的 Schwarzschild 时空来算一算 Komar 积分。归一化为 $n_{\mu}n^{\mu}=-1$ 与 $\sigma_{\mu}\sigma^{\mu}=+1$ 的法矢量，其非零分量为
\begin{equation*}
n_{0} = -\left(1-\frac{2GM}{r}\right)^{1/2}, \qquad \sigma_{1} = \left(1-\frac{2GM}{r}\right)^{-1/2},
\tag{6.39}
\end{equation*}
其余分量均为零。于是我们有
\begin{equation*}
n_{\mu}\sigma_{\nu}\nabla^{\mu}K^{\nu} = -\nabla^{0}K^{1}.
\tag{6.40}
\end{equation*}
Killing 矢量为 $K^{\mu}=(1,0,0,0)$，因此我们可以毫不费力地算出
\begin{align*}
\nabla^{0}K^{1} &= g^{00}\nabla_{0}K^{1}\\
&= g^{00}\left(\partial_{0}K^{1}+\Gamma^{1}_{0\lambda}K^{\lambda}\right)\\
&= g^{00}\Gamma^{1}_{00}K^{0}\\
&= -\left(1-\frac{2GM}{r}\right)^{-1}\frac{GM}{r^{2}}\left(1-\frac{2GM}{r}\right)\\
&= -\frac{GM}{r^{2}}.
\tag{6.41}
\end{align*}
无穷远处二维球面上的度规是
\begin{equation*}
\gamma^{(2)}_{ij}\,\mathrm{d}x^{i}\mathrm{d}x^{j} = r^{2}\left(\mathrm{d}\theta^{2}+\sin^{2}\theta\,\mathrm{d}\phi^{2}\right),
\tag{6.42}
\end{equation*}
于是
\begin{equation*}
\sqrt{\gamma^{(2)}} = r^{2}\sin\theta.
\tag{6.43}
\end{equation*}

% ---- 书页 252 / p267 ----
把这一切合在一起，Schwarzschild 黑洞的能量就是
\begin{align*}
E_{R} &= \frac{1}{4\pi G}\int d\theta\, d\phi\, r^{2}\sin\theta\left(\frac{GM}{r^{2}}\right)\\
&= M.
\tag{6.44}
\end{align*}
这当然正是我们想要的结果，同时也解释了式 (6.35) 中所选取的归一化。

尽管得到了正确的答案，我们还是应该想一想刚才到底发生了什么。具体来说，我们是把流 $J^{\mu}_{R}=K_{\nu}R^{\mu\nu}$ 在一个类空切片上积分而得到这个能量的，并且发现结果可以写成空间无穷远处的积分。但对 Schwarzschild 来说，度规满足真空 Einstein 方程 $R_{\mu\nu}=0$；这样一来，要从对 $J^{\mu}_{R}$ 的积分中得到非零的结果似乎就很困难，正如式 (6.29) 那样。可是如果想一想最大延拓 Schwarzschild 解的结构，我们就会意识到，可以画出两类类空切片：一类穿过虫洞一直延伸到第二个渐近区域，另一类则终止于奇点。如果切片穿过虫洞，另一个渐近区域就会为 $\partial\Sigma$ 提供另一个组成部分，从而对式 (6.38) 提供另一份贡献；这份贡献会恰好抵消，于是总能量确实会是零。如果切片撞上了奇点，我们就不知道该怎么处理了。尽管如此，无论哪种情形，把我们的结果 (6.44) 当作正确的答案都是合理的。关键在于，式 (6.38) 只涉及空间无穷远处的贡献，所以无论内部发生什么，它都应当是能量的一个有效表达式。我们甚至可以设想内部出现随时间变化的行为；只要 $K^{\mu}$ \emph{渐近地}是类时 Killing 矢量，Komar 能量就是良定义的。比如说，我们可以考虑从一颗初始静态的恒星出发的球对称引力坍缩。直接在 $\Sigma$ 上计算积分 (6.35)，会给出一个合理的总质量答案，而且当恒星坍缩为黑洞时，这个总质量不应改变（我们设想的是球对称，因此引力辐射无法把能量带到无穷远）。于是，在坍缩之前有效的 Komar 积分 (6.38)，在坍缩成黑洞之后也尽可以放心地解释为能量。当然，出于某些目的，我们或许想容许通过引力辐射损失能量，那种情形下就必须小心对待如何把切片延伸到无穷远；人们可以在未来类光无限远处定义一个“Bondi 质量”（Bondi mass），用它来追踪通过辐射损失的能量。

关于 Komar 公式的另一个疑虑是，它是否真的是我们应当认作“能量”的那种东西——能量通常是指与时间平移不变性相联系的守恒量。支持这种解释的最好理由很简单：$E_{R}$ 无疑是某种守恒量，而且它与我们认为应当是 Schwarzschild 能量的数值相一致（在 Newton 极限下对一团质量合集也是如此，你可以自行验证），那么它还能是什么呢？另一种做法是考虑广义相对论的哈密顿表述，在渐近平直时空中仔细定义时间平移的生成元，再把它认同为总能量。这项工作最先由 Arnowitt、Deser 和 Misner 完成，他们的结果就是所谓的\textbf{ADM 能量}（ADM energy）。在一个
% ---- 书页 253 / p268 ----
渐近平直时空中，我们可以像在微扰论中那样把度规写成
\begin{equation*}
g_{\mu\nu} = \eta_{\mu\nu} + h_{\mu\nu},
\tag{6.45}
\end{equation*}
只不过这里我们只要求分量 $h_{\mu\nu}$ 在空间无穷远处很小，而不必处处都小。于是 ADM 能量可以写成空间无穷远处一个二维球面上的积分：
\begin{equation*}
E_{\mathrm{ADM}} = \frac{1}{16\pi G}\int_{\partial\Sigma} d^{2}x\,\sqrt{\gamma^{(2)}}\,\sigma^{i}\left(\partial_{j}h^{j}{}_{i}-\partial_{i}h^{j}{}_{j}\right),
\tag{6.46}
\end{equation*}
其中空间指标用 $\delta^{ij}$（无穷远处的空间度规）升起。这个公式看上去依赖于坐标的选取，但在我们的假设之下实际上是良定义的。如果 $h_{\mu\nu}$ 在无穷远处不依赖于时间，可以验证 ADM 能量与 Komar 能量实际上是一致的。这让我们更加相信 Komar 积分确实表示能量。不过，在某种意义上 ADM 能量更为“体面”；例如，如果存在与引力非最小耦合的长程标量场，Komar 积分就可能遇到麻烦。但就我们眼前的目的而言，Komar 能量已经完全够用了。

我们希望某个被叫作“能量”的东西具有这样一种品质：对任何物理位形它都是正的；否则，一个零能量的态就可以衰变成一些正能量的碎片和一些负能量的碎片。第 4 章讨论过的能量条件为物质场给出了正能量的概念，但我们可能会担心，引力的负能量贡献会带来麻烦。所幸在 GR 中我们有\textbf{正能量定理}（positive energy theorem），它最先由 Schoen 和 Yau（译注：原文排印作 “Shoen”，应为 Schoen。）证明：

无奇点的、满足 Einstein 方程和主能量条件的渐近平直时空，其 ADM 能量是非负的。而且，Minkowski 时空是唯一使 ADM 能量为零的此类时空。

如果容许奇点存在，反例显然是有的，比如 $M<0$ 的 Schwarzschild 解。不过，如果一个带奇点的时空（比如 $M>0$ 的 Schwarzschild 解）是作为无奇点初数据的演化结果而达到的，定理就仍然适用。这样看来，在广义相对论中我们似乎是安全的，不会遇到负能量的孤立系统。

最后，我们可以转向自旋（角动量）；讨论完能量之后，这一步轻而易举。设想我们有一个转动 Killing 矢量 $R=\partial_{\phi}$。与时间平移的情形完全类似，我们可以定义一个守恒流
\begin{equation*}
J^{\mu}_{\phi} = R_{\nu}R^{\mu\nu},
\tag{6.47}
\end{equation*}
由它将得到把守恒角动量 $J$ 表示成空间无穷远处积分的表达式，

% ---- 书页 254 / p269 ----
\begin{equation*}
\boxed{\,J = -\frac{1}{8\pi G}\int_{\partial\Sigma} d^{2}x\,\sqrt{\gamma^{(2)}}\, n_{\mu}\sigma_{\nu}\nabla^{\mu}R^{\nu}.\,}
\tag{6.48}
\end{equation*}
（“J”既被用来表示流又被用来表示角动量，“R”既是转动 Killing 矢量又是 Ricci 张量，这都挺遗憾。可是字母就只有这么多，只能将就了。）与能量一样，即使 $R^{\mu}$ 只是渐近地才是 Killing 矢量，这个表达式仍然有效。注意，这里的归一化与能量积分中的不同；例如，可以通过对弱引力场中缓慢运动的物体求出该表达式的值，来为这个归一化提供依据。

\section{带电（Reissner–Nordström）黑洞}

现在我们转向表示带电黑洞的精确解。这类解与现实的天体物理情形关系并不太大；在真实世界中，一个高电荷黑洞会通过与黑洞附近物质的相互作用迅速被中和。不过，带电黑洞仍然展示了更一般情形的若干重要特征。在这个问题中，完整的球对称性依然存在；因此我们知道可以把度规写成
\begin{equation*}
\mathrm{d}s^{2} = -e^{2\alpha(r,t)}\mathrm{d}t^{2} + e^{2\beta(r,t)}\mathrm{d}r^{2} + r^{2}\mathrm{d}\Omega^{2}.
\tag{6.49}
\end{equation*}
然而现在我们不再处于真空，因为黑洞会带有非零的电磁场，而电磁场又充当能量-动量的源。电磁场的能量-动量张量为
\begin{equation*}
T_{\mu\nu} = F_{\mu\rho}F_{\nu}{}^{\rho} - \frac{1}{4}g_{\mu\nu}F_{\rho\sigma}F^{\rho\sigma},
\tag{6.50}
\end{equation*}
其中 $F_{\mu\nu}$ 是电磁场强张量。既然问题具有球对称性，最一般的场强张量就只能有如下分量
\begin{gather*}
F_{tr} = f(r,t) = -F_{rt}\\
F_{\theta\phi} = g(r,t)\sin\theta = -F_{\phi\theta},
\tag{6.51}
\end{gather*}
其中 $f(r,t)$ 和 $g(r,t)$ 是两个有待场方程确定的函数，没有写出的分量均为零。$F_{tr}$ 对应于径向电场，而 $F_{\theta\phi}$ 对应于径向磁场。如果你正纳闷那个 $\sin\theta$ 是从哪儿来的，回忆一下：不应当依赖于 $\theta$ 和 $\phi$ 的量，是磁场的径向分量 $B^{r}=\epsilon^{01\mu\nu}F_{\mu\nu}$。对球对称的度规，
\begin{equation*}
\epsilon^{\rho\sigma\mu\nu} = \frac{1}{\sqrt{-g}}\tilde{\epsilon}^{\rho\sigma\mu\nu}
\end{equation*}
与 $(\sin\theta)^{-1}$ 成正比，所以我们希望在 $F_{\theta\phi}$ 中带有一个 $\sin\theta$ 因子。

% ---- 书页 255 / p270 ----
这种情形下的场方程既有 Einstein 方程，也有麦克斯韦方程组：
\begin{gather*}
g^{\mu\nu}\nabla_{\mu}F_{\nu\sigma} = 0\\
\nabla_{[\mu}F_{\nu\rho]} = 0.
\tag{6.52}
\end{gather*}
这两组方程彼此耦合在一起，因为电磁场强张量通过能量-动量张量进入 Einstein 方程，而度规又显式地出现在麦克斯韦方程组之中。

然而，这些困难并非不可克服；与真空情形中我们如法炮制的步骤相类似的做法，同样会给出带电情形的解。我们不会显式地写出各个步骤，只把最后的结果抄下来。这个解就是\textbf{Reissner–Nordström（RN）度规}（Reissner–Nordström metric），
\begin{equation*}
\boxed{\,\mathrm{d}s^{2} = -\Delta\,\mathrm{d}t^{2} + \Delta^{-1}\mathrm{d}r^{2} + r^{2}\mathrm{d}\Omega^{2},\,}
\tag{6.53}
\end{equation*}
其中
\begin{equation*}
\Delta = 1 - \frac{2GM}{r} + \frac{G(Q^{2}+P^{2})}{r^{2}}.
\tag{6.54}
\end{equation*}
在这个表达式中，$M$ 依旧解释为黑洞的质量；$Q$ 是总电荷，$P$ 是总磁荷。孤立的磁荷（磁单极子）从未在自然界中被观测到，但这并不妨碍我们写下它们如果真的存在时将会产生的度规。\footnote{本章采用的单位制中，库仑定律里没有 $4\pi$ 因子。要想与其他章节比较，把每一处出现的 $Q$ 或 $P$ 除以 $4\pi$ 即可。}有很好的理论理由让我们相信，如果各种相互作用在极高能量下“大统一”，磁单极子就可能存在，但它们必定非常重，而且极为稀少。当然，即使不存在任何磁单极子，黑洞也可能带上磁荷。实际上，电荷与磁荷进入度规的方式完全相同，因此在表达式中保留 $P$ 并不会引入任何额外的复杂之处。保守的读者尽可以随意令 $P=0$。与这个解相联系的电磁场为
\begin{gather*}
E_{r} = F_{rt} = \frac{Q}{r^{2}}\\
B_{r} = \frac{F_{\theta\phi}}{r^{2}\sin\theta} = \frac{P}{r^{2}}.
\tag{6.55}
\end{gather*}
这些场的 $1/r^{2}$ 依赖关系正是我们在平直空间中所熟悉的；当然，我们知道这依赖于我们对径向坐标的精确选取。RN 度规在 $r=0$ 处有一个名副其实的曲率奇点，这可以通过计算曲率不变标量 $R_{\mu\nu\rho\sigma}R^{\mu\nu\rho\sigma}$ 来验证。视界结构
% ---- 书页 256 / p271 ----
却比 Schwarzschild 情形复杂得多。在上面关于事件视界的讨论中我们曾提出，如果我们聪明地选取了坐标，使得这个条件在 $r$ 的某个固定值处得到满足，那么 $g^{rr}=0$ 就会是定位事件视界的一个有用判据。幸运的是，式 (6.53) 的坐标恰好具有这种性质，于是事件视界将位于
\begin{equation*}
g^{rr}(r) = \Delta(r) = 1 - \frac{2GM}{r} + \frac{G(Q^{2}+P^{2})}{r^{2}} = 0.
\tag{6.56}
\end{equation*}
这将发生在
\begin{equation*}
r_{\pm} = GM \pm \sqrt{G^{2}M^{2} - G(Q^{2}+P^{2})}.
\tag{6.57}
\end{equation*}
%FIG{6.2}: {Reissner–Nordström 解的函数 $\Delta(r)=1-2GM/r+G(Q^{2}+P^{2})/r^{2}$；零点标示事件视界所在的位置。}
如图 6.2 所示，随着 $GM^{2}$ 与 $Q^{2}+P^{2}$ 相对大小的不同，这可能给出两个、一个或零个解。于是我们逐个考察各种情形。

\subsection*{情形一：$GM^{2}<Q^{2}+P^{2}$}

这种情形下，系数 $\Delta$ 总为正（永不为零），度规在 $(t,r,\theta,\phi)$ 坐标中一直直到 $r=0$ 都完全正则。坐标 $t$ 总是类时的，而 $r$ 总是类空的。但 $r=0$ 处依然有奇点，只是现在它是一条类时线。由于不存在事件视界，观者前往奇点、看过一眼再返回报告所见，没有任何阻碍。这就是前面讨论过的裸奇点。对测地线的仔细分析表明，奇点是排斥的——类时测地线永远不会与 $r=0$ 相交；它们只是靠近，然后掉头离去。（类光测地线可以到达奇点，非测地线的类时曲线也可以。）

% ---- 书页 257 / p272 ----
%FIG{6.3}: {Reissner–Nordström 解在 $Q^{2}+P^{2}>GM^{2}$ 时的共形图。原点处有一个裸奇点。}

当 $r\to\infty$ 时解趋近于平直时空，而且正如我们刚刚看到的，因果结构看上去处处正常。因此共形图就会与 Minkowski 空间的一模一样，只不过现在 $r=0$ 是一个奇点，如图 6.3 所示。

奇点的裸露冒犯了我们的体面感，也同样冒犯了宇宙监督假设。实际上，我们决不应指望引力坍缩的结果会是一个 $GM^{2}<Q^{2}+P^{2}$ 的黑洞。粗略地说，这个条件说的是，黑洞的总能量比单独电磁场贡献的能量还要少——也就是说，携带电荷的那些物质，其质量必须为负。因此这个解通常被认为是非物理的。还要注意，这个时空中不存在任何 Cauchy 面，因为类时线可以在奇点处开始或终止。

\subsection*{情形二：$GM^{2}>Q^{2}+P^{2}$}

我们预期这种情形适用于现实的引力坍缩；电磁场中的能量小于总能量。此时度规系数 $\Delta(r)$ 在 $r$ 很大和很小时都为正，而在两个零点 $r_{\pm}=GM\pm\sqrt{G^{2}M^{2}-G(Q^{2}+P^{2})}$ 之间为负。度规在 $r_{+}$ 和 $r_{-}$ 两处都有坐标奇点；与 Schwarzschild 的情形一样，这两处都可以通过坐标变换消去。

由 $r=r_{\pm}$ 定义的曲面都是类光的，而且二者都是事件视界。$r=0$ 处的奇点是一条类时线，而不像 Schwarzschild 中那样是一个类空曲面。如果你是一个从远处落入黑洞的观者，$r_{+}$ 就好比 Schwarzschild 度规中的 $2GM$：在这个半径处，$r$ 从类空坐标切换成类时坐标，你必然朝 $r$ 减小的方向运动。黑洞外面的目击者也会看到同样的
% ---- 书页 258 / p273 ----
现象，与不带电黑洞外面的情形相同——落入的观者看上去运动得越来越慢，红移也越来越厉害。

但是，从 $r_{+}$ 一路坠向越来越小的半径这种不可避免的坠落，只会持续到你抵达类光曲面 $r=r_{-}$ 为止；在那里 $r$ 重新变回类空坐标，朝 $r$ 减小方向的运动便可以止住。因此你并不一定要撞上 $r=0$ 处的奇点；这是意料之中的，因为 $r=0$ 是一条类时线（因而未必处在你的未来之中）。实际上，你既可以选取继续走向 $r=0$，也可以开始朝 $r$ 增大的方向折返，重新穿过 $r=r_{-}$ 处的类光曲面。此后 $r$ 将再一次成为类时坐标，只不过取向反了过来；你被迫朝 $r$ \emph{增大}的方向运动。最终你会再一次被抛过 $r=r_{+}$，就仿佛从一个白洞里冒出来、进入宇宙的其余部分。从这里你又可以选取回到黑洞中去——这一次是另一个洞，与你最初进入的那个不同——并且随你喜欢把这样的旅行重复多少次。这个小故事对应于图 6.4 中的共形图；当然，通过选取适当的坐标并把 Reissner–Nordström 度规解析延拓到极致，可以更严格地把它推导出来。

%FIG{6.4}: {Reissner–Nordström 解在 $GM^{2}>Q^{2}+P^{2}$ 时的共形图。黑洞外部区域有无穷多个副本。}

% ---- 书页 259 / p274 ----
%FIG{6.5}: {极端 Reissner–Nordström 解（$GM^{2}=Q^{2}+P^{2}$）的共形图。原点处有一个裸奇点，外部区域则有无穷多个。}

这番情景里有多少是科学、有多少是科学幻想呢？大概科学不多。如果你从一位身处黑洞内部、正要穿过 $r=r_{-}$ 处事件视界的观者的角度来考虑这个世界，你会注意到，这位观者可以回溯时间，看到外部（渐近平直）宇宙的整部历史——至少从黑洞看来是如此。但是，观者是在自己有限的固有时间内看到这段（无限长的）历史的——于是，当观者接近 $r_{-}$ 时，到达他们的任何信号都会被无限蓝移。因此，任何进入 RN 黑洞的非球对称微扰，都很可能会剧烈地扰动我们所描述的几何。真实的几何会是什么样子很难说，但没有很好的理由相信它必定包含无穷多个通过种种虫洞彼此连通的渐近平直区域。\footnote{关于这个问题的一些研究，见 E. Poisson and W. Israel, \emph{Phys.\ Rev.\ D} \textbf{41}, 1796 (1990).}

\subsection*{情形三：$GM^{2}=Q^{2}+P^{2}$}

这种情形称为\textbf{极端}（extremal）Reissner–Nordström 解。一方面，极端黑洞是一件有趣的理论玩具；在研究黑洞在量子引力中扮演的角色时，人们经常考察这个解。在超对称理论中，极端黑洞可以让某些对称性保持不破缺，这给计算带来极大的方便。另一方面，它看上去是不稳定的，因为只要加入一点点物质，它就会变成情形二。

极端黑洞的 $\Delta(r)$ 只在单个半径 $r=GM$ 处为零。这代表一个事件视界，但 $r$ 坐标永远不是类时的；它在 $r=GM$ 处变成类光的，而在两侧都是类空的。与其他情形一样，$r=0$ 处的奇点是一条类时线。所以对这个黑洞来说，你同样可以避开奇点，继续向未来运动到渐近平直区域的额外副本中去，只不过奇点总是“在左边”。共形图如图 6.5 所示。

极端黑洞的一个迷人性质是，质量在某种意义上被电荷平衡了。更具体地说，两个携带同号电荷的极端黑洞会彼此产生引力吸引，却又彼此产生电磁排斥，而这两种效应恰好抵消。事实上，对于相耦合的 Einstein–Maxwell 方程组，我们能找到\emph{精确}解，表示处在稳定位形中的任意数目的此类黑洞。为看清这一点，先回到 Reissner–Nordström 度规本身；为简单起见，我们只考虑电荷而不考虑磁荷。在极端性条件下 $GM^{2}=Q^{2}$，度规化为
\begin{equation*}
\mathrm{d}s^{2} = -\left(1-\frac{GM}{r}\right)^{2}\mathrm{d}t^{2} + \left(1-\frac{GM}{r}\right)^{-2}\mathrm{d}r^{2} + r^{2}\mathrm{d}\Omega^{2}.
\tag{6.58}
\end{equation*}
通过定义一个平移的径向坐标
\begin{equation*}
\rho = r - GM,
\tag{6.59}
\end{equation*}

% ---- 书页 260 / p275 ----
度规便取各向同性的形式
\begin{equation*}
\mathrm{d}s^{2} = -H^{-2}(\rho)\mathrm{d}t^{2} + H^{2}(\rho)\left[\mathrm{d}\rho^{2} + \rho^{2}\mathrm{d}\Omega^{2}\right],
\tag{6.60}
\end{equation*}
其中
\begin{equation*}
H(\rho) = 1 + \frac{GM}{\rho}.
\tag{6.61}
\end{equation*}
因为 $\mathrm{d}\rho^{2}+\rho^{2}\mathrm{d}\Omega^{2}$ 正是三个空间维度中的平直度规，我们同样可以把式 (6.60) 写成
\begin{equation*}
\mathrm{d}s^{2} = -H^{-2}(\vec{x})\mathrm{d}t^{2} + H^{2}(\vec{x})\left[\mathrm{d}x^{2} + \mathrm{d}y^{2} + \mathrm{d}z^{2}\right],
\tag{6.62}
\end{equation*}
其中 $H$ 可以写成
\begin{equation*}
H = 1 + \frac{GM}{|\vec{x}|}.
\tag{6.63}
\end{equation*}
在原先的 $r$ 坐标中，极端解的电场可以用矢势 $A_{\mu}$ 表示为
\begin{equation*}
E_{r} = F_{rt} = \frac{Q}{r^{2}} = \partial_{r}A_{0},
\tag{6.64}
\end{equation*}
其中矢势的类时分量为
\begin{equation*}
A_{0} = -\frac{Q}{r},
\tag{6.65}
\end{equation*}
并且我们设想空间分量为零（已经把磁场设为零）。在新的 $\rho$ 坐标下，再加上极端性条件 $Q^{2}=GM^{2}$，上式变为
\begin{equation*}
A_{0} = -\frac{\sqrt{G}M}{\rho+GM},
\tag{6.66}
\end{equation*}
或者等价地写成
\begin{equation*}
\sqrt{G}\,A_{0} = H^{-1}-1.
\tag{6.67}
\end{equation*}
但现在让我们忘掉 $H$ 满足式 (6.61) 这件事，径直把度规 (6.62) 和静电势 (6.67) 代入 Einstein 方程与麦克斯韦方程组，设想 $H$ 不依赖于时间（$\partial_{0}H=0$）但除此之外不受任何约束。我们可以直截了当地证明（见习题）：任何一个不依赖于时间、且满足
\begin{equation*}
\nabla^{2}H = 0,
\tag{6.68}
\end{equation*}
（其中 $\nabla^{2}=\partial_{x}^{2}+\partial_{y}^{2}+\partial_{z}^{2}$）的函数 $H(\vec{x})$，都能让这两组方程同时得到满足。这其实就是拉普拉斯方程，而且很容易写出所有在无穷远处表现良好的解；它们的形式为

% ---- 书页 261 / p276 ----
\begin{equation*}
H = 1 + \sum_{a=1}^{N}\frac{GM_{a}}{|\vec{x}-\vec{x}_{a}|},
\tag{6.69}
\end{equation*}
其中 $\vec{x}_{a}$ 给出某组 $N$ 个空间点。这些点描述了 $N$ 个极端 RN 黑洞所在的位置，各黑洞的质量为 $M_{a}$，电荷为 $Q_{a}=\sqrt{G}M_{a}$。这个多极端黑洞度规无疑是 Einstein 方程最引人注目的精确解之一。

\section{转动（Kerr）黑洞}

关于带电解，我们本来还可以多谈不少细节，但我们还是接着讲转动黑洞吧。在这种情形下寻找度规的精确解要困难得多，因为我们已经放弃了球对称性。我们转而寻找绕转动轴具有轴对称性、同时又是稳态的（即存在一个类时 Killing 矢量的）解。虽然 Schwarzschild 解和 Reissner–Nordström 解在广义相对论问世后不久就被发现了，但转动黑洞的解直到 1963 年才由 Kerr 找到。他的结果——\textbf{Kerr 度规}（Kerr metric）——就是下面这团乱麻：
\begin{equation*}
\boxed{\begin{aligned}
\mathrm{d}s^{2} &= -\left(1-\frac{2GMr}{\rho^{2}}\right)\mathrm{d}t^{2} - \frac{2GMar\sin^{2}\theta}{\rho^{2}}\left(\mathrm{d}t\,\mathrm{d}\phi + \mathrm{d}\phi\,\mathrm{d}t\right)\\
&\quad + \frac{\rho^{2}}{\Delta}\mathrm{d}r^{2} + \rho^{2}\mathrm{d}\theta^{2} + \frac{\sin^{2}\theta}{\rho^{2}}\left[(r^{2}+a^{2})^{2} - a^{2}\Delta\sin^{2}\theta\right]\mathrm{d}\phi^{2},
\end{aligned}}
\tag{6.70}
\end{equation*}
其中
\begin{equation*}
\boxed{\,\Delta(r) = r^{2} - 2GMr + a^{2}\,}
\tag{6.71}
\end{equation*}
以及
\begin{equation*}
\boxed{\,\rho^{2}(r,\theta) = r^{2} + a^{2}\cos^{2}\theta.\,}
\tag{6.72}
\end{equation*}
两个常数 $M$ 和 $a$ 参数化了各种可能的解。要验证质量 $M$ 确实等于 Komar 能量 (6.38) 并不难，只是繁琐；而 $a$ 是单位质量的角动量，
\begin{equation*}
a = J/M,
\tag{6.73}
\end{equation*}
其中 $J$ 是 Komar 角动量 (6.48)。要把电荷和磁荷 $Q$ 与 $P$ 也包括进来很容易，只需把 $2GMr$ 换成 $2GMr-G(Q^{2}+$

% ---- 书页 262 / p277 ----
$P^{2})$；得到的结果就是\textbf{Kerr–Newman 度规}（Kerr–Newman metric）。与之相联系的 1-形式势具有如下非零分量
\begin{equation*}
A_{t} = \frac{Qr - Pa\cos\theta}{\rho^{2}}, \qquad A_{\phi} = \frac{-Qar\sin^{2}\theta + P(r^{2}+a^{2})\cos\theta}{\rho^{2}}.
\tag{6.74}
\end{equation*}
所有本质性的现象在电荷不存在时依然保持，所以从现在起我们令 $Q=P=0$。

坐标 $(t,r,\theta,\phi)$ 就是所谓的\textbf{Boyer–Lindquist 坐标}（Boyer–Lindquist coordinates）。很容易验证，当 $a\to 0$ 时它们退化为 Schwarzschild 坐标。然而，如果保持 $a$ 不动而让 $M\to 0$，我们得到的仍是平直时空，只不过不是用普通的极坐标来表示的。此时度规变为
\begin{equation*}
\mathrm{d}s^{2} = -\mathrm{d}t^{2} + \frac{r^{2}+a^{2}\cos^{2}\theta}{r^{2}+a^{2}}\mathrm{d}r^{2} + \left(r^{2}+a^{2}\cos^{2}\theta\right)^{2}\mathrm{d}\theta^{2} + \left(r^{2}+a^{2}\right)\sin^{2}\theta\,\mathrm{d}\phi^{2},
\tag{6.75}
\end{equation*}
我们认出它的空间部分正是用椭球坐标表示的平直空间，如图 6.6 所示。这些坐标与欧几里得三维空间中的笛卡尔坐标之间的关系是
\begin{align*}
x &= \left(r^{2}+a^{2}\right)^{1/2}\sin\theta\cos\phi\\
y &= \left(r^{2}+a^{2}\right)^{1/2}\sin\theta\sin\phi\\
z &= r\cos\theta.
\tag{6.76}
\end{align*}

度规 (6.70) 有两个 Killing 矢量，二者都是显而易见的；由于度规系数不依赖于 $t$ 和 $\phi$，$K=\partial_{t}$ 与 $R=\partial_{\phi}$ 都是 Killing 矢量。当然，$R^{\mu}$ 表达的是解的轴对称性。矢量 $K^{\mu}$ 并不与 $t=$ 常数的超曲面正交，实际上也不与任何超曲面正交；因此这个度规是稳态的，但不是
% 本块末句在原书书页263顶部继续（“static. This makes sense; ...”），下一块请用 %JOIN% 续接本句。
